\documentclass[12pt,a4paper]{article}

\usepackage[utf8]{inputenc}
\usepackage{amsmath}
\usepackage{graphicx}
\usepackage{float}
\usepackage{hyperref}
\usepackage{geometry}
\geometry{margin=1in}

\begin{document}

\begin{titlepage}
    \centering
    \vspace*{6.5cm}
    {\LARGE \textbf{Network Intrusion Detection System} \\[0.5em]
        \textbf{Powered by Machine Learning}}\\[2cm]
    {\large Phase 2: Data Understanding }\\[2cm]
    {\large Amira BOUAFIF \& Hajer JEMAA}\\[0.5cm]
    \vfill
    {\large 2025-2026}\\[1cm]
\end{titlepage}

\section{Introduction}

This report presents the Data Understanding phase of the CRISP-DM methodology applied
to the \textbf{CIC-IDS2017} dataset. The objective of this phase is to gain a clear
understanding of the dataset, including its structure and characteristics, and to
identify any potential issues that will be addressed in the Data Preparation phase.

\section{Dataset Overview}

The \textbf{CIC-IDS2017 dataset} consists of network traffic data collected over five
working days (Monday to Friday). The raw traffic was captured and then transformed into
tabular data containing flow-level statistics (e.g., duration, packet counts, and rates)
using CICFlowMeter. Comprising over 2.8 million rows and 79 columns, the dataset
includes both benign and malicious traffic across 14 different attack types, such as
DoS, Brute Force, DDoS, and SQL injection.

\section{Methodology}

The CRISP-DM (Cross-Industry Standard Process for Data Mining) methodology is followed
for the Data Understanding phase. CRISP-DM is a widely used framework for structuring
machine learning projects. It consists of six phases:

\begin{itemize}
    \item Business Understanding
    \item Data Understanding
    \item Data Preparation
    \item Modeling
    \item Evaluation
    \item Deployment
\end{itemize}

The Data Understanding phase involves loading the data, describing its structure and
content, exploring its characteristics, and finally verifying its quality.

\section{Data Loading}

The dataset is distributed across eight CSV files stored in a shared data directory.
Each file is assigned a short alias for use throughout the analysis, as presented in
Table~\ref{tab:file_aliases}.

\begin{table}[H]
    \centering
    \caption{Original filenames and their assigned aliases.}
    \label{tab:file_aliases}
    \begin{tabular}{|l|l|}
        \hline
        \textbf{Original filename}                                            & \textbf{Alias}     \\
        \hline
        \texttt{Monday-WorkingHours.pcap\_ISCX.csv}                           & Monday             \\
        \hline
        \texttt{Tuesday-WorkingHours.pcap\_ISCX.csv}                          & Tuesday            \\
        \hline
        \texttt{Wednesday-workingHours.pcap\_ISCX.csv}                        & Wednesday          \\
        \hline
        \texttt{Thursday-WorkingHours-Morning-WebAttacks.pcap\_ISCX.csv}      & Thursday-AM        \\
        \hline
        \texttt{Thursday-WorkingHours-Afternoon-Infilteration.pcap\_ISCX.csv} & Thursday-PM        \\
        \hline
        \texttt{Friday-WorkingHours-Morning.pcap\_ISCX.csv}                   & Friday-AM          \\
        \hline
        \texttt{Friday-WorkingHours-Afternoon-DDos.pcap\_ISCX.csv}            & Friday-PM-DDoS     \\
        \hline
        \texttt{Friday-WorkingHours-Afternoon-PortScan.pcap\_ISCX.csv}        & Friday-PM-PortScan \\
        \hline
    \end{tabular}
\end{table}

As shown in Table~\ref{tab:file_aliases}, each file is named after the day and
session it corresponds to. The files are loaded into memory using pandas with the
\texttt{low\_memory=False} option to ensure consistent type inference across the full
column range. All eight files are ingested successfully into a dictionary of DataFrames
indexed by alias, which serves as the primary data structure for all subsequent analyses.

\section{Data Description}

This section examines the surface properties of the CIC-IDS2017 dataset in the order
they are addressed in the analysis: file dimensions and column consistency, preliminary
formatting corrections, label inventory, data combination, data types, duplicate columns,
and finally the semantic interpretation of the features.

\subsection{File Shapes and Column Consistency}

The two following subsections verify that the eight files share a common structure,
which is a prerequisite for safe concatenation and for using a single representative
file in subsequent analyses.

\subsubsection{File Shapes}

The dataset comprises eight CSV files generated by CICFlowMeter from five days of
network traffic captures. Table~\ref{tab:file_shapes} presents the dimensions
of each file after loading.

\begin{table}[H]
    \centering
    \caption{Dimensions of each file after loading.}
    \label{tab:file_shapes}
    \begin{tabular}{|l|r|r|}
        \hline
        \textbf{File alias} & \textbf{Rows}        & \textbf{Columns} \\
        \hline
        Monday              & 529\,918             & 79               \\
        \hline
        Tuesday             & 445\,909             & 79               \\
        \hline
        Wednesday           & 692\,703             & 79               \\
        \hline
        Thursday-AM         & 170\,366             & 79               \\
        \hline
        Thursday-PM         & 288\,602             & 79               \\
        \hline
        Friday-AM           & 191\,033             & 79               \\
        \hline
        Friday-PM-DDoS      & 225\,745             & 79               \\
        \hline
        Friday-PM-PortScan  & 286\,467             & 79               \\
        \hline
        \textbf{Total}      & \textbf{2\,830\,743} & 79               \\
        \hline
    \end{tabular}
\end{table}

\subsubsection{Column Consistency Check}
\label{sec:colconsistency}

As shown in Table~\ref{tab:file_shapes}, each file contains exactly 79 columns.
A programmatic check confirms that all eight files share the same 79 column names
in the same order, using the Monday file as the reference. This uniformity is an
expected consequence of a single CICFlowMeter pipeline being applied to all captures
and guarantees that a single file can serve as a representative sample for the
schema-level analyses that follow, and that vertical concatenation into a unified
DataFrame can be performed without column alignment issues.

\subsection{Preliminary Formatting Corrections}

Before any analysis is performed, two formatting defects are identified and corrected.
Both corrections are purely cosmetic and do not alter any numerical feature value.

\subsubsection{Whitespace in Column Names}

Inspection of the raw CSV headers reveals that 65 out of 79 column names contain
leading or trailing whitespace characters (for example, \texttt{` Destination Port'}
in place of \texttt{`Destination Port'}). This is a documented artefact of
CICFlowMeter's CSV export format. Left uncorrected, such whitespace causes
\texttt{KeyError} exceptions when columns are accessed by name in pandas
(e.g., \texttt{df['Label']} would fail). All column names are stripped across all
eight files before any further step is performed.

\subsubsection{Label Encoding Corruption}

Three label values in the Thursday-AM file are found to contain mojibake characters:
the hyphen separating ``Web Attack'' from the attack sub-type has been corrupted into
the Unicode replacement character \texttt{U+FFFD}. Table~\ref{tab:label_fixes}
presents the affected labels alongside their corrected counterparts.

\begin{table}[H]
    \centering
    \caption{Label encoding corrections applied to the Thursday-AM file.}
    \label{tab:label_fixes}
    \begin{tabular}{|l|l|}
        \hline
        \textbf{Corrupted label}            & \textbf{Corrected label}            \\
        \hline
        \texttt{Web Attack ? Brute Force}   & \texttt{Web Attack - Brute Force}   \\
        \hline
        \texttt{Web Attack ? XSS}           & \texttt{Web Attack - XSS}           \\
        \hline
        \texttt{Web Attack ? Sql Injection} & \texttt{Web Attack - Sql Injection} \\
        \hline
    \end{tabular}
\end{table}

As shown in Table~\ref{tab:label_fixes}, the three affected labels are exclusively
Web Attack variants confined to the Thursday-AM file. The correction is applied via
a direct string replacement before any label-based grouping or filtering is performed.

\subsection{Label Inventory}

A label inventory is compiled per file in order to verify that pandas has not
introduced unexpected types and to document which attack classes are present in
each capture session. The result confirms that the \texttt{Label} column is of
type \texttt{object} across all files and that it is the only non-numeric column
in the dataset. Table~\ref{tab:label_inventory} presents the classes present
in each file.

\begin{table}[H]
    \centering
    \caption{Label inventory per file.}
    \label{tab:label_inventory}
    \begin{tabular}{|l|l|}
        \hline
        \textbf{File alias} & \textbf{Classes present}                           \\
        \hline
        Monday              & BENIGN                                             \\
        \hline
        Tuesday             & BENIGN, FTP-Patator, SSH-Patator                   \\
        \hline
        Wednesday           & BENIGN, DoS GoldenEye, DoS Hulk, DoS Slowhttptest, \\
                            & DoS slowloris, Heartbleed                          \\
        \hline
        Thursday-AM         & BENIGN, Web Attack - Brute Force,                  \\
                            & Web Attack - XSS, Web Attack - Sql Injection       \\
        \hline
        Thursday-PM         & BENIGN, Infiltration                               \\
        \hline
        Friday-AM           & BENIGN, Bot                                        \\
        \hline
        Friday-PM-DDoS      & BENIGN, DDoS                                       \\
        \hline
        Friday-PM-PortScan  & BENIGN, PortScan                                   \\
        \hline
    \end{tabular}
\end{table}

As illustrated in Table~\ref{tab:label_inventory}, Monday contains exclusively
benign traffic and serves as the baseline capture session. Each attack type appears
in exactly one file.

\subsection{Data Combination and Preview}

Given the verified column consistency established in Section~\ref{sec:colconsistency},
the eight files are concatenated vertically into a single unified DataFrame
\texttt{df\_all} of \textbf{2,830,743 rows and 79 columns}. All subsequent analyses
in this report operate on this combined dataset unless stated otherwise.

\subsection{Data Types}

The type composition of the dataset is assessed using the Monday file as a
representative sample. Table~\ref{tab:dtypes} presents the distribution of
column types across the 79 columns.

\begin{table}[H]
    \centering
    \caption{Data type distribution across the 79 columns.}
    \label{tab:dtypes}
    \begin{tabular}{|l|r|l|}
        \hline
        \textbf{Type}    & \textbf{Count} & \textbf{Description}                                           \\
        \hline
        \texttt{int64}   & 54             & Integer-valued flow counters and TCP flag counts.              \\
        \hline
        \texttt{float64} & 24             & Continuous flow statistics: rates, means, standard deviations. \\
        \hline
        \texttt{object}  & 1              & The \texttt{Label} column, the sole non-numeric field.         \\
        \hline
        \textbf{Total}   & \textbf{79}    &                                                                \\
        \hline
    \end{tabular}
\end{table}

As shown in Table~\ref{tab:dtypes}, all 78 feature columns are in numeric form
and require no categorical encoding at this stage. The absence of any additional
\texttt{object} columns confirms that no typing errors were introduced by pandas
during loading. It is worth noting that the two rate features \texttt{Flow Bytes/s}
and \texttt{Flow Packets/s}, despite being typed as \texttt{float64}, are
conceptually throughput quantities. Their floating-point typing is a consequence
of the presence of \texttt{Inf} values in those columns, which are observed when
\texttt{Flow Duration} equals zero. This observation is documented here and
examined further in the Data Quality Verification section.

\subsection{Duplicate Column}

A column named \texttt{Fwd Header Length.1} is found to be an exact duplicate of
\texttt{Fwd Header Length} across all 2\,830\,743 rows. The suffix \texttt{.1} is
appended automatically by pandas when a CSV file contains two identically named
columns, a known artefact of early CICFlowMeter versions. This redundant column
carries no additional information and will be removed during the Data Preparation phase.

\subsection{Feature Description}

The two following subsections document the semantic interpretation of the 78 numeric
features and address a conceptual aspect of the dataset that directly influences the
modelling strategy.

\subsubsection{Feature Groups}

All 78 numeric features are produced by CICFlowMeter from raw packet captures.
Each row represents a single bidirectional network flow, defined as the complete
set of packets exchanged between two endpoints sharing the same 5-tuple (source IP,
destination IP, source port, destination port, and protocol) within a session window.
The features are statistical aggregates computed over all packets belonging to that
flow. Table~\ref{tab:feature_groups} presents the feature groups and their semantics.

\begin{table}[H]
    \centering
    \caption{Feature groups in CIC-IDS2017 and their semantics.}
    \label{tab:feature_groups}
    \begin{tabular}{|p{3cm}|p{4.5cm}|p{6.5cm}|}
        \hline
        \textbf{Group} & \textbf{Representative features}                                                         & \textbf{Captured information} \\
        \hline
        Flow metadata
                       & \texttt{Flow Duration}, \texttt{Protocol}
                       & Total session duration and transport protocol (TCP, UDP, ICMP).                                                          \\
        \hline
        Packet counts
                       & \texttt{Total Fwd Packets}, \texttt{Total Backward Packets}
                       & Number of packets transmitted in each direction.                                                                         \\
        \hline
        Packet lengths
                       & \texttt{Fwd Packet Length Max/Min/Mean/Std}
                       & Statistical distribution of packet sizes in bytes.                                                                       \\
        \hline
        Inter-arrival times
                       & \texttt{Flow IAT Mean}, \texttt{Fwd IAT Total}
                       & Time elapsed between consecutive packets; captures the temporal rhythm of traffic.                                       \\
        \hline
        TCP flags
                       & \texttt{SYN Flag Count}, \texttt{ACK Flag Count}, \texttt{FIN Flag Count}
                       & Frequency of TCP control flags; characterises connection establishment and teardown.                                     \\
        \hline
        Flow rates
                       & \texttt{Flow Bytes/s}, \texttt{Flow Packets/s}
                       & Throughput metrics. Subject to undefined values when \texttt{Flow Duration} equals zero.                                 \\
        \hline
        Header lengths
                       & \texttt{Fwd Header Length}, \texttt{Bwd Header Length}
                       & Protocol header overhead per direction, in bytes.                                                                        \\
        \hline
        Window size
                       & \texttt{Init\_Win\_bytes\_forward /backward}
                       & Initial TCP receive window advertised at connection establishment.                                                       \\
        \hline
    \end{tabular}
\end{table}  

\newpage                                        
\begin{table}[H]
    \centering
    \begin{tabular}{|p{3cm}|p{4.5cm}|p{6.5cm}|}
        \hline
        Active / Idle times
                       & \texttt{Active Mean}, \texttt{Idle Mean}
                       & Duration of active transmission periods and idle intervals within the flow.\\
        \hline
        Bulk transfer
                       & \texttt{Fwd Avg Bytes/Bulk}, \texttt{Bwd Avg Bulk Rate}
                       & Bulk transfer statistics. All values are zero in this dataset due to a known
        CICFlowMeter limitation.                                                                                                                  \\
        \hline
        Target
                       & \texttt{Label}
                       & Class label: \texttt{BENIGN} or the name of the attack type.                                                             \\
        \hline
    \end{tabular}
\end{table}

As illustrated in Table~\ref{tab:feature_groups}, the feature set covers a broad
spectrum of network flow characteristics. The notation \textbf{Fwd} refers to the
client-to-server direction, while \textbf{Bwd} refers to the server-to-client direction.

\subsubsection{Representation of Multi-Event Attacks at the Flow Level}

A distinction must be drawn between the operational definition of certain attack types
and the granularity at which the dataset represents them. Some attacks produce a
detectable anomaly within a single network interaction: an SQL injection may manifest
in the payload of a single HTTP request/response exchange. However, several attack
families are inherently multi-event phenomena whose defining characteristic is a
pattern that emerges across many flows, not within any single one. Three such cases
are present in CIC-IDS2017:

\begin{itemize}
    \item \textbf{DDoS}: a large volume of connection attempts directed at a single
          victim within a short time window. No individual flow is necessarily abnormal;
          the threat resides in the volume and concurrence across flows.
    \item \textbf{PortScan}: a systematic sequence of connection attempts across many
          destination ports. A single probe flow may appear entirely benign; the scan
          pattern only emerges when many flows are observed together.
    \item \textbf{Brute Force (FTP-Patator, SSH-Patator)}: repeated authentication
          attempts across multiple sessions. The attack is characterised by its rate
          and repetition across flows, not by any single authentication exchange.
\end{itemize}

Because CIC-IDS2017 is produced in a controlled experimental environment where the
execution schedule of each attack is known, all flows captured during an attack window
are labelled with the corresponding class. A flow-level classifier therefore learns
per-flow statistical fingerprints correlated with each attack class, such as
characteristic packet rates, flag distributions, or inter-arrival time patterns that
emerge at the individual flow level during an attack campaign. The full operational
definition of multi-event attacks, which requires aggregating evidence across many
flows over time, is not directly captured at this granularity. In a production context,
per-flow model scores would consequently be consumed by an alerting layer applying
temporal aggregation and thresholding to produce actionable incident alerts.

\section{Data Exploration}

This section explores the distribution of classes and the pairwise correlations between
features. All analyses are performed on the raw, uncleaned dataset; findings are
therefore preliminary and will be revisited after the cleaning steps of the Data
Preparation phase.

\subsection{Class Distribution}

The following subsections examine the distribution of the target variable at two levels:
the global distribution across the entire combined dataset, and the per-file breakdown
by capture session.

\subsubsection{Global Class Distribution}

table~\ref{tab:class_distribution} presents the sample count, percentage, and
learnability rating for each of the 15 classes present in the combined dataset.
The learnability rating is defined as follows: \textit{Sufficient} for classes with at
least 10,000 samples, \textit{Low} for classes between 1,000 and 9,999 samples,
\textit{Insufficient} for classes between 100 and 999 samples, and \textit{Critical}
for classes with fewer than 100 samples.

\begin{table}[H]
    \centering
    \caption{Global class distribution with learnability assessment.}
    \label{tab:class_distribution}
    \begin{tabular}{|l|r|r|l|}
        \hline
        \textbf{Class}             & \textbf{Count}     & \textbf{Percentage} & \textbf{Learnability} \\
        \hline
        BENIGN                     & 2,273,097          & 80.300\%            & Sufficient*           \\
        \hline
        DoS Hulk                   & 231,073            & 8.163\%             & Sufficient            \\
        \hline
        PortScan                   & 158,930            & 5.614\%             & Sufficient            \\
        \hline
        DDoS                       & 128,027            & 4.523\%             & Sufficient            \\
        \hline
        DoS GoldenEye              & 10,293             & 0.364\%             & Sufficient            \\
        \hline
        FTP-Patator                & 7,938              & 0.280\%             & Low                   \\
        \hline
        SSH-Patator                & 5,897              & 0.208\%             & Low                   \\
        \hline
        DoS slowloris              & 5,796              & 0.205\%             & Low                   \\
        \hline
        DoS Slowhttptest           & 5,499              & 0.194\%             & Low                   \\
        \hline
        Bot                        & 1,966              & 0.069\%             & Low                   \\
        \hline
        Web Attack - Brute Force   & 1,507              & 0.053\%             & Low                   \\
        \hline
        Web Attack - XSS           & 652                & 0.023\%             & Insufficient          \\
        \hline
        Infiltration               & 36                 & 0.001\%             & Critical              \\
        \hline
        Web Attack - Sql Injection & 21                 & 0.001\%             & Critical              \\
        \hline
        Heartbleed                 & 11                 & 0.000\%             & Critical              \\
        \hline
        \textbf{Total}             & \textbf{2,830,743} & \textbf{100\%}      &                       \\
        \hline
        \multicolumn{4}{|l|}{\small * BENIGN has sufficient samples but its 80.3\% dominance
        creates class imbalance.}                                                                     \\
        \hline
    \end{tabular}
\end{table}

As shown in Table~\ref{tab:class_distribution}, the dataset presents a severe class
imbalance: BENIGN traffic accounts for 80.3\% of all rows while the 14 attack classes
combined represent only 19.7\%, yielding an imbalance ratio of 4.1 benign samples for
every attack sample. Attack class sizes span four orders of magnitude, from 231,073
samples for DoS Hulk down to 11 for Heartbleed. Three classes fall in the Critical
category, namely Infiltration, Web Attack - Sql Injection, and Heartbleed, for which
standard supervised learning is unlikely to generalise reliably. La
Figure~\ref{fig:class_global} presents the full distribution on both linear and
logarithmic scales, and la Figure~\ref{fig:attack_only} focuses exclusively on the
attack classes to make the imbalance across minor classes visible.

\begin{figure}[H]
    \centering
    \includegraphics[width=\textwidth]{../figs/class_distribution_global.png}
    \caption{Global class distribution across all 15 classes (linear and log scales).}
    \label{fig:class_global}
\end{figure}

\begin{figure}[H]
    \centering
    \includegraphics[width=\textwidth]{../figs/Attack_only_Distribution.png}
    \caption{Attack-only class distribution with learnability colour coding.}
    \label{fig:attack_only}
\end{figure}

\subsubsection{Per-File Class Distribution}

La Figure~\ref{fig:per_file_dist} presents the class distribution broken down by
capture file, and la Figure~\ref{fig:benign_vs_attack} presents the BENIGN-to-attack
ratio per file.

\begin{figure}[H]
    \centering
    \includegraphics[width=\textwidth]{../figs/class_distribution_per_file.png}
    \caption{Class distribution per capture file.}
    \label{fig:per_file_dist}
\end{figure}

\begin{figure}[H]
    \centering
    \includegraphics[width=\textwidth]{../figs/benign_vs_attack_per_file.png}
    \caption{BENIGN vs attack ratio per file.}
    \label{fig:benign_vs_attack}
\end{figure}

As illustrated in Figure~\ref{fig:per_file_dist} and Figure~\ref{fig:benign_vs_attack},
each attack type is confined to exactly one capture file. FTP-Patator and SSH-Patator
appear exclusively in Tuesday, Web Attack variants exclusively in Thursday-AM, and DDoS
exclusively in Friday-PM-DDoS. This structural property has a direct consequence for
evaluation strategy: a file-based train/test split would entirely exclude certain attack
classes from training, making the model unable to recognise them at evaluation time.
The BENIGN-to-attack ratio also varies substantially across files, from 100\% benign on
Monday to approximately equal proportions in certain afternoon sessions, a heterogeneity
that may introduce source bias if not handled carefully during splitting. The implications
for split strategy selection are examined in Section~\ref{sec:quality}.

\subsection{Feature Correlation}

A pairwise Pearson correlation matrix is computed on the 78 numeric features of the
combined dataset. Table~\ref{tab:high_corr} presents a selection of the 39 feature
pairs whose absolute correlation exceeds the threshold of 0.95.

\begin{table}[H]
    \centering
    \caption{Selection of highly correlated feature pairs (absolute correlation above 0.95).}
    \label{tab:high_corr}
    \begin{tabular}{|l|l|r|}
        \hline
        \textbf{Feature A}          & \textbf{Feature B}   & \textbf{Correlation} \\
        \hline
        Bwd Packet Length Mean      & Avg Bwd Segment Size & 1.000000             \\
        \hline
        Fwd PSH Flags               & SYN Flag Count       & 1.000000             \\
        \hline
        Total Fwd Packets           & Subflow Fwd Packets  & 1.000000             \\
        \hline
        Fwd Packet Length Mean      & Avg Fwd Segment Size & 1.000000             \\
        \hline
        Total Backward Packets      & Subflow Bwd Packets  & 1.000000             \\
        \hline
        Fwd Header Length           & Fwd Header Length.1  & 1.000000             \\
        \hline
        Total Length of Bwd Packets & Subflow Bwd Bytes    & 1.000000             \\
        \hline
        Total Length of Fwd Packets & Subflow Fwd Bytes    & 0.999999             \\
        \hline
        Flow Duration               & Fwd IAT Total        & 0.998560             \\
        \hline
        Packet Length Mean          & Average Packet Size  & 0.997790             \\
        \hline
        Flow Packets/s              & Fwd Packets/s        & 0.987562             \\
        \hline
        \multicolumn{3}{|l|}{\small ... 28 additional pairs above 0.95 (full list available
        in the notebook).}                                                        \\
        \hline
    \end{tabular}
\end{table}

As shown in Table~\ref{tab:high_corr}, many of the perfect correlations reflect
algebraic relationships within CICFlowMeter's feature computation: for instance,
\texttt{Avg Fwd Segment Size} is by definition equal to \texttt{Fwd Packet Length Mean},
and subflow packet counts are identical to total packet counts under standard CICFlowMeter
settings. These redundant features carry no additional information and are candidates for
removal during the feature selection step of the Data Preparation phase. La
Figure~\ref{fig:corr_heatmap} presents the correlation heatmap restricted to the
features involved in at least one high-correlation pair.

\begin{figure}[H]
    \centering
    \includegraphics[width=\textwidth]{../figs/correlation_heatmap.png}
    \caption{Pearson correlation heatmap for the subset of features involved in
        high-correlation pairs.}
    \label{fig:corr_heatmap}
\end{figure}

\section{Data Quality Verification}
\label{sec:quality}

This section assesses the quality of the raw combined dataset across seven dimensions:
missing and infinite values, duplicate rows, constant and near-constant features, feature
distribution properties, outlier quantification, feature-target association, per-file
distribution drift, and finally the data leakage risk and split strategy selection. All
checks are performed on the uncleaned data; the findings presented here define the
cleaning agenda for the Data Preparation phase.

\subsection{Missing and Infinite Values}

The two following subsections document the column-level extent of missing and infinite
values and assess their impact at the class level to determine whether dropping affected
rows would introduce survivorship bias against rare attack classes.

\subsubsection{Column-Level Analysis}

Missing values (NaN) and infinite values (Inf) are examined separately. The NaN check
reveals that \textbf{1,358 rows} contain at least one missing value, representing
0.0480\% of the dataset. The affected column is exclusively \texttt{Flow Bytes/s},
whose missing values arise from a division by zero when \texttt{Flow Duration} equals
zero: CICFlowMeter attempts to compute flow rate as bytes divided by duration, producing
an undefined result for instantaneous flows. The Inf check reveals that \textbf{2,867
    rows} contain at least one infinite value (0.1013\% of the dataset), affecting two
columns: \texttt{Flow Packets/s} (2,867 Inf entries) and \texttt{Flow Bytes/s} (1,509
Inf entries), both arising from the same root cause. Table~\ref{tab:nan_inf} summarises
these findings.

\begin{table}[H]
    \centering
    \caption{Summary of NaN and Inf values per affected column.}
    \label{tab:nan_inf}
    \begin{tabular}{|l|r|r|r|}
        \hline
        \textbf{Column}         & \textbf{NaN Count} & \textbf{Inf Count} & \textbf{\% of rows}             \\
        \hline
        \texttt{Flow Bytes/s}   & 1,358              & 1,509              & 0.0480\% (NaN) / 0.0533\% (Inf) \\
        \hline
        \texttt{Flow Packets/s} & 0                  & 2,867              & 0.1013\%                        \\
        \hline
    \end{tabular}
\end{table}

As shown in Table~\ref{tab:nan_inf}, the combined row-level impact of NaN and Inf
values is small in absolute terms, with fewer than 0.15\% of all rows affected.

\subsubsection{Class-Level Impact}

Although the overall proportion of affected rows is small, it is necessary to verify
that NaN and Inf values are not concentrated in rare attack classes, in which case
dropping these rows could disproportionately reduce the already limited sample size of
those classes. La Figure~\ref{fig:nan_inf_impact} presents the class-level impact of
NaN and Inf rows.

\begin{figure}[H]
    \centering
    \includegraphics[width=\textwidth]{../figs/class_impact_nan_inf.png}
    \caption{Class-level impact of NaN and Inf rows in absolute count and as a
        percentage of each class total.}
    \label{fig:nan_inf_impact}
\end{figure}

As illustrated in Figure~\ref{fig:nan_inf_impact}, the three Critical classes
(Heartbleed, Web Attack - Sql Injection, Infiltration) and all Insufficient classes are
entirely unaffected, with zero NaN or Inf rows. The class with the highest relative
impact is Bot, which loses 10 rows out of 1,966 (0.51\%), a negligible proportion.
Dropping rows containing NaN or Inf values is therefore a safe strategy that introduces
no survivorship bias against rare classes.

\subsection{Duplicate Rows}

Duplicate row analysis is conducted in two stages: an exact duplicate check, which
identifies rows that are identical in all 79 columns including the label, and a
contradictory duplicate check, which identifies rows whose feature vectors are identical
but whose labels differ.

\subsubsection{Exact Duplicates}

The exact duplicate check reveals that \textbf{308,381 rows} are exact duplicates of
an earlier row, representing 10.8940\% of the combined dataset. Table~\ref{tab:dupes}
presents the distribution of duplicate rows by class.

\begin{table}[H]
    \centering
    \caption{Exact duplicate rows by class.}
    \label{tab:dupes}
    \begin{tabular}{|l|r|}
        \hline
        \textbf{Class}           & \textbf{Duplicate Rows} \\
        \hline
        BENIGN                   & 176,613                 \\
        \hline
        PortScan                 & 68,111                  \\
        \hline
        DoS Hulk                 & 58,224                  \\
        \hline
        SSH-Patator              & 2,678                   \\
        \hline
        FTP-Patator              & 2,005                   \\
        \hline
        DoS slowloris            & 411                     \\
        \hline
        DoS Slowhttptest         & 271                     \\
        \hline
        Web Attack - Brute Force & 37                      \\
        \hline
        Bot                      & 13                      \\
        \hline
        DDoS                     & 11                      \\
        \hline
        DoS GoldenEye            & 7                       \\
        \hline
        \textbf{Total}           & \textbf{308,381}        \\
        \hline
    \end{tabular}
\end{table}

As shown in Table~\ref{tab:dupes}, duplicate rows are common in network traffic
datasets because many flows share identical statistical profiles, for instance short
single-packet DNS queries. Exact duplicates in both features and label can however
inflate accuracy metrics during evaluation by placing near-identical flows in both the
training and test partitions. The decision on whether to deduplicate and at what
granularity is deferred to the Data Preparation phase.

\subsubsection{Contradictory Duplicates}
\label{sec:contradictory}

A contradictory duplicate is defined as a group of rows that share an identical feature
vector (verified via hash-based comparison) but are assigned different class labels. The
analysis identifies \textbf{698 such groups} involving a total of \textbf{7,020 rows}.
Table~\ref{tab:contradictory} presents the classes involved and their respective row
counts within contradictory groups.

\begin{table}[H]
    \centering
    \caption{Classes involved in contradictory duplicate groups.}
    \label{tab:contradictory}
    \begin{tabular}{|l|r|}
        \hline
        \textbf{Class} & \textbf{Rows in conflicting groups} \\
        \hline
        DoS Hulk       & 5,224                               \\
        \hline
        PortScan       & 1,053                               \\
        \hline
        BENIGN         & 736                                 \\
        \hline
        DDoS           & 6                                   \\
        \hline
        DoS slowloris  & 1                                   \\
        \hline
        \textbf{Total} & \textbf{7,020}                      \\
        \hline
    \end{tabular}
\end{table}

As shown in Table~\ref{tab:contradictory}, all classes involved in contradictions are
well-represented (above 1,000 samples), and no Critical or Insufficient class is affected.
The origin of these contradictions lies in the flow-aggregation process: CICFlowMeter
computes statistical summaries over all packets in a flow, and two flows from different
sessions (one benign, one malicious) can produce statistically identical feature vectors
if their traffic patterns are indistinguishable at the aggregation level. This constitutes
an irreducible source of label noise. No classifier can correctly label all rows in a
contradictory group from features alone, and accuracy metrics will therefore have a hard
ceiling for the affected classes. All rows belonging to contradictory groups will be
dropped during Data Preparation, as they carry no recoverable signal.

\subsection{Constant and Near-Constant Features}

The following subsections identify features that carry no discriminative information due
to an absence of variation across the entire dataset, distinguishing between features
that are strictly constant and those that are dominated by a single value.

\subsubsection{Strictly Constant Features}
\label{sec:constant}

Eight numeric columns take a single unique value (zero) across all 2,830,743 rows.
These columns are listed in Table~\ref{tab:constant}.

\begin{table}[H]
    \centering
    \caption{Strictly constant features and their unique value.}
    \label{tab:constant}
    \begin{tabular}{|l|r|}
        \hline
        \textbf{Feature}              & \textbf{Unique value} \\
        \hline
        \texttt{Bwd PSH Flags}        & 0                     \\
        \hline
        \texttt{Bwd URG Flags}        & 0                     \\
        \hline
        \texttt{Fwd Avg Bytes/Bulk}   & 0                     \\
        \hline
        \texttt{Fwd Avg Packets/Bulk} & 0                     \\
        \hline
        \texttt{Fwd Avg Bulk Rate}    & 0                     \\
        \hline
        \texttt{Bwd Avg Bytes/Bulk}   & 0                     \\
        \hline
        \texttt{Bwd Avg Packets/Bulk} & 0                     \\
        \hline
        \texttt{Bwd Avg Bulk Rate}    & 0                     \\
        \hline
    \end{tabular}
\end{table}

As shown in Table~\ref{tab:constant}, all eight constant columns are bulk transfer
statistics. These features are intended to measure large continuous data transfers where
the sender pushes data in bulk without waiting for acknowledgement between segments. In
practice, a known bug in CICFlowMeter causes these values to be output as zero regardless
of the actual traffic, a limitation documented in the research community for this specific
tool and dataset version. These features carry zero information and will be removed
unconditionally during Data Preparation.

\subsubsection{Near-Constant Features}

Four additional features are found to be near-constant, defined as having a single
dominant value in at least 99.5\% of all rows. Table~\ref{tab:nearconstant} presents
these features along with their dominance percentages.

\begin{table}[H]
    \centering
    \caption{Near-constant features with their dominant value and dominance rate.}
    \label{tab:nearconstant}
    \begin{tabular}{|l|r|r|r|}
        \hline
        \textbf{Feature}        & \textbf{Dominant value} & \textbf{Dominance (\%)} & \textbf{Unique values} \\
        \hline
        \texttt{Fwd URG Flags}  & 0                       & 99.989\%                & 2                      \\
        \hline
        \texttt{CWE Flag Count} & 0                       & 99.989\%                & 2                      \\
        \hline
        \texttt{RST Flag Count} & 0                       & 99.976\%                & 2                      \\
        \hline
        \texttt{ECE Flag Count} & 0                       & 99.976\%                & 2                      \\
        \hline
    \end{tabular}
\end{table}

As illustrated in Table~\ref{tab:nearconstant}, the non-zero values for all four
near-constant features appear exclusively in BENIGN traffic (315 rows for
\texttt{Fwd URG Flags} and \texttt{CWE Flag Count}, 686 rows for \texttt{RST Flag Count},
and 689 rows for \texttt{ECE Flag Count}), meaning these features provide no signal
for distinguishing any attack class from benign traffic. They are strong candidates for
removal during feature selection in the Data Preparation phase.

\subsection{Feature Distribution Properties}

This section characterises the statistical distribution of the 78 numeric features across
four dimensions: skewness, visual separation between benign and attack flows, zero
inflation, and scale heterogeneity.

\subsubsection{Skewness}

Among the 78 numeric features, \textbf{41 exhibit a skewness with absolute value greater
    than 5}, indicating heavily right-tailed distributions. This is characteristic of network
traffic data, where the majority of flows are short and low-volume while a small fraction
involves very large or very long exchanges. La Figure~\ref{fig:skewness} presents the
skewness values for all 78 features, with those exceeding the threshold highlighted in red.

\begin{figure}[H]
    \centering
    \includegraphics[width=\textwidth]{../figs/feature_skewness.png}
    \caption{Skewness of all 78 numeric features; features with absolute skewness above
        5 are highlighted in red (41 features).}
    \label{fig:skewness}
\end{figure}

It is noted that skewness does not invalidate tree-based models such as XGBoost, which
are invariant to monotonic transformations of the input. However, highly skewed features
may reduce interpretability and pose challenges for distance-based or linear models.
Logarithmic or power transforms will be evaluated empirically during the Data Preparation
phase.

\subsubsection{Feature Distributions by Class}

La Figure~\ref{fig:feat_dist} presents overlaid histograms for eight key features,
contrasting the distributions of BENIGN and ATTACK flows on the combined dataset.
Values are clipped to the [q1, q99] range for visualisation clarity.

\begin{figure}[H]
    \centering
    \includegraphics[width=\textwidth]{../figs/feature_distribution_benign_vs_attack.png}
    \caption{Distributions of eight key features for BENIGN versus ATTACK flows,
        clipped to the [q1, q99] range.}
    \label{fig:feat_dist}
\end{figure}

As illustrated in Figure~\ref{fig:feat_dist}, \texttt{Flow Bytes/s} and
\texttt{Flow Packets/s} exhibit distinct peaks for the two classes, suggesting partial
separability in those dimensions. Forward and backward packet length means show
different profiles for attack flows, particularly for DoS-type attacks that generate
many small packets. The inter-arrival time mean (\texttt{Flow IAT Mean}) also differs
between classes, as attacks with rapid packet bursts produce significantly shorter
inter-arrival times than benign traffic. These observations are based on binary grouping
and serve as a preliminary indication of feature utility; multi-class patterns will be
examined during the modelling phase.



\subsubsection{Feature Scale Heterogeneity}

Feature value ranges span vastly different magnitudes across the 78 numeric columns,
from binary-like features taking values in \{0, 1\} to flow rates and byte counts
exceeding hundreds of millions. This heterogeneity is inherent to the CICFlowMeter
feature set, which mixes packet counts, durations expressed in microseconds, byte
sizes, and dimensionless ratios. Notably, \texttt{Fwd Header Length} and its duplicate
\texttt{Fwd Header Length.1} exhibit anomalous negative minimum values
($-$32,212,234,632), which are likely integer overflow artefacts in CICFlowMeter and
will be investigated during Data Preparation. For the planned XGBoost baseline model,
feature scaling is not required as tree-based models are invariant to the scale of
input features. Scaling would however be necessary if distance-based or linear models
are explored in later phases.

\subsection{Outlier Analysis}

An outlier is defined, for each feature, as any value falling below the 1st percentile
or above the 99th percentile of that feature's distribution (computed after replacing
Inf values with NaN). Table~\ref{tab:outliers} presents the top 10 features by total
outlier count.

\begin{table}[H]
    \centering
    \caption{Top 10 features by outlier count outside the [q1, q99] range.}
    \label{tab:outliers}
    \begin{tabular}{|l|r|r|r|}
        \hline
        \textbf{Feature}             & \textbf{Below q1} & \textbf{Above q99} & \textbf{Total (\%)} \\
        \hline
        \texttt{Fwd Packets/s}       & 28,308            & 27,618             & 55,926 (1.98\%)     \\
        \hline
        \texttt{Destination Port}    & 17,568            & 28,302             & 45,870 (1.62\%)     \\
        \hline
        \texttt{Flow IAT Mean}       & 3,478             & 28,186             & 31,664 (1.12\%)     \\
        \hline
        \texttt{Flow Duration}       & 2,982             & 28,308             & 31,290 (1.11\%)     \\
        \hline
        \texttt{Flow IAT Min}        & 2,891             & 28,308             & 31,199 (1.10\%)     \\
        \hline
        \texttt{Flow IAT Max}        & 2,981             & 28,140             & 31,121 (1.10\%)     \\
        \hline
        \texttt{Fwd Header Length.1} & 1,809             & 26,845             & 28,654 (1.01\%)     \\
        \hline
        \texttt{Fwd Header Length}   & 1,809             & 26,845             & 28,654 (1.01\%)     \\
        \hline
        \texttt{Flow Packets/s}      & 28,279            & 334                & 28,613 (1.01\%)     \\
        \hline
        \texttt{Fwd IAT Min}         & 17                & 28,306             & 28,323 (1.00\%)     \\
        \hline
    \end{tabular}
\end{table}

As shown in Table~\ref{tab:outliers}, the class-level breakdown of outliers reveals
that for \texttt{Destination Port}, FTP-Patator contributes 7,937 outlier rows out of
its total 7,938 samples, indicating that extreme port values are a defining
characteristic of this attack rather than noise. By contrast, for features such as
\texttt{Flow IAT Mean} and \texttt{Flow Duration}, BENIGN traffic accounts for the
majority of outlier rows, reflecting the natural long-tailed behaviour of normal traffic.
Outlier handling during Data Preparation must therefore be informed by this class
breakdown: outliers characteristic of a specific attack class should be preserved, while
those attributable to normal traffic variability may be candidates for clipping or
winsorisation.

\subsection{Feature-Target Association}

Mutual information is computed between each of the 78 features and a binary target
(BENIGN vs ATTACK) on a stratified subsample of 50,000 rows, drawn from the 2,827,876
rows that remain after removing NaN and Inf values. La Figure~\ref{fig:mi_top20}
presents the 20 features with the highest mutual information scores.

\begin{figure}[H]
    \centering
    \includegraphics[width=\textwidth]{../figs/feature_mi_top20.png}
    \caption{Top 20 features by mutual information score with respect to the binary
        BENIGN vs ATTACK target.}
    \label{fig:mi_top20}
\end{figure}

As illustrated in Figure~\ref{fig:mi_top20}, packet length statistics emerge as the
most discriminative features: \texttt{Packet Length Variance} and \texttt{Packet Length
    Std} both score 0.3364, followed by \texttt{Average Packet Size} at 0.3307. These
results are consistent with domain knowledge, as attacks such as DoS and DDoS alter
the packet size distribution measurably at the flow level. At the opposite extreme,
all bulk transfer features score at or near zero (below 0.0002), consistent with the
finding that these features are always zero due to the CICFlowMeter limitation
documented in Section~\ref{sec:constant}. It is noted that these scores reflect
association with a binary target only; feature importance rankings may differ for the
multi-class classification task.

\subsection{Per-File Feature Distribution Drift}

To isolate the effect of temporal and environmental variation from the presence of
attacks, the median values of six representative features are compared across the eight
capture files using BENIGN-only traffic. Any observed shift is attributable exclusively
to session-specific factors such as time of day, network load, or traffic composition.

The analysis reveals substantial drift for certain features: \texttt{Flow Duration}
exhibits a maximum-to-minimum median ratio of \textbf{560.09} across files, flagged
as high drift, driven by the distinctly shorter flows observed in the Thursday-PM
session (median of 306 microseconds vs. over 171,000 in Friday-PM-DDoS). By contrast,
\texttt{Total Fwd Packets} remains stable across all files with a median of 2 in every
session, indicating that this feature is robust to session-level variation. Features
with high drift ratios may cause the model to learn session-specific patterns rather
than genuine attack signatures, and will be monitored during cross-validation in the
modelling phase.

\subsection{Data Leakage Risk and Split Strategy}

The following two aspects of data leakage are examined before selecting a split
strategy. First, the presence of non-target object-type features is checked: no
such columns exist beyond the \texttt{Label} column itself, confirming that no
string-valued identifier or metadata column can leak class information into the feature
space. Second, contradictory duplicate groups are confirmed to be present (698 groups,
as documented in Section~\ref{sec:contradictory}), introducing an irreducible Bayes
error for any classifier.

Table~\ref{tab:split_strategy} presents the trade-offs between five candidate
splitting strategies for CIC-IDS2017.

\begin{table}[H]
    \centering
    \caption{Comparison of candidate train/test split strategies for CIC-IDS2017.}
    \label{tab:split_strategy}
    \begin{tabular}{|l|l|l|l|}
        \hline
        \textbf{Strategy}          & \textbf{All classes?} & \textbf{Leakage risk} & \textbf{Complexity} \\
        \hline
        Random                     & Not guaranteed        & High                  & Low                 \\
        \hline
        File-based (by day)        & No                    & Low                   & Low                 \\
        \hline
        Stratified split           & Yes                   & Medium                & Easy                \\
        \hline
        Stratified K-Fold          & Yes                   & Medium                & Medium              \\
        \hline
        Temporal + stratified/file & Yes                   & Low                   & Hard                \\
        \hline
    \end{tabular}
\end{table}

As shown in Table~\ref{tab:split_strategy}, two strategies are immediately excluded
by structural constraints. The random split does not guarantee that all 14 attack types
appear in both training and test partitions, which is particularly problematic for
Critical classes such as Heartbleed (11 samples), and risks placing temporally adjacent
flows from the same session in both partitions, inflating performance metrics. The
file-based split fails because each attack type is confined to exactly one file: any
file excluded from training renders an entire attack class unseen by the model.
Stratified K-Fold is therefore selected as the recommended primary evaluation strategy,
as it guarantees proportional representation of every class in each fold while providing
multiple evaluation runs for more robust metric estimation. Implementation is deferred
to the Data Preparation phase.

\section{Conclusion}

This report presented the Data Understanding phase of the CRISP-DM methodology applied
to the CIC-IDS2017 dataset. The analysis was conducted on the raw, uncleaned data and
covered four successive tasks: data loading, data description, data exploration, and
data quality verification.

The dataset was found to comprise approximately 2.83 million flow-level records across
79 columns, representing 15 traffic classes generated over five working days. The
examination of the class distribution revealed a severe imbalance, with BENIGN traffic
accounting for 80.3\% of all samples, while three attack classes (Heartbleed,
Infiltration, and Web Attack - Sql Injection) each contained fewer than 100 samples,
placing them in the Critical learnability category. The most consequential structural
property identified was that each attack type is confined to exactly one capture file,
a constraint that eliminated file-based splitting as a viable evaluation strategy and
motivated the adoption of Stratified K-Fold cross-validation.

The data quality verification identified several issues that must be addressed before
modelling. NaN and Inf values were found to be concentrated in two rate features
(\texttt{Flow Bytes/s} and \texttt{Flow Packets/s}), affecting fewer than 0.15\% of
rows with no bias against rare classes. Exact duplicate rows represented 10.89\% of
the dataset, and 698 contradictory duplicate groups were identified, introducing an
irreducible source of label noise. Eight features were found to be strictly constant
across all rows due to a known CICFlowMeter bug, and four additional features were
near-constant, both groups being candidates for unconditional removal. Feature analysis
further revealed pervasive right-skewness in 41 out of 78 features, 39 highly correlated
feature pairs reflecting algebraic redundancies in the CICFlowMeter feature set, and
substantial per-file distribution drift in BENIGN traffic confirming the presence of
temporal variation across capture sessions.

It is noted that all statistics computed in this phase, including correlation
coefficients, skewness values, and descriptive summaries, were computed on raw data
containing Inf values and are therefore preliminary. All findings will be recomputed
on cleaned data following the Data Preparation phase.

The Data Preparation phase will address the identified issues in the following order:
replacement of Inf values with NaN followed by row removal, elimination of contradictory
and exact duplicate rows, removal of constant and near-constant features, resolution of
highly correlated feature pairs through feature selection, label encoding, implementation
of Stratified K-Fold splitting, treatment of class imbalance, and normalisation of
numerical features where required by the chosen model family.

\end{document}
